\documentclass[12pt,a4paper]{article}

% --- CẤU HÌNH TIẾNG VIỆT HOÀN HẢO CHO OVERLEAF (PDFLATEX) ---
\usepackage[T5]{fontenc} % Bắt buộc để Overleaf hiển thị đúng font tiếng Việt với PdfLaTeX
\usepackage[utf8]{inputenc}
\usepackage[vietnam]{babel}
\usepackage{times} % Sử dụng font Times New Roman tiêu chuẩn học thuật

\usepackage{amsmath}
\usepackage{amsfonts}
\usepackage{amssymb}
\usepackage{graphicx}
\usepackage{hyperref}
\usepackage{listings}
\usepackage{color}
\usepackage{float}
\usepackage{geometry}
\usepackage{tikz}
\usetikzlibrary{shapes,arrows,positioning,calc}

\geometry{
 a4paper,
 total={170mm,257mm},
 left=30mm,
 right=20mm,
 top=20mm,
 bottom=20mm
}

% Cấu hình hiển thị mã nguồn (C++ & CMake)
\definecolor{codegreen}{rgb}{0,0.6,0}
\definecolor{codegray}{rgb}{0.5,0.5,0.5}
\definecolor{codepurple}{rgb}{0.58,0,0.82}
\definecolor{backcolour}{rgb}{0.95,0.95,0.92}

\lstdefinestyle{mystyle}{
    backgroundcolor=\color{backcolour},   
    commentstyle=\color{codegreen},
    keywordstyle=\color{blue},
    numberstyle=\tiny\color{codegray},
    stringstyle=\color{codepurple},
    basicstyle=\ttfamily\footnotesize,
    breakatwhitespace=false,         
    breaklines=true,                 
    captionpos=b,                    
    keepspaces=true,                 
    numbers=left,                    
    numbersep=5pt,                  
    showspaces=false,                
    showstringspaces=false,
    showtabs=false,                  
    tabsize=2
}
\lstset{style=mystyle}

% Cấu hình liên kết hyperref
\hypersetup{
    colorlinks=true,
    linkcolor=blue,
    filecolor=magenta,      
    urlcolor=cyan,
    pdftitle={Báo cáo kỹ thuật: Hệ thống Thang máy thông minh nhận diện khuôn mặt AI},
    pdfpagemode=FullScreen,
}

\begin{document}

% =======================================================
% TRANG BÌA HỌC THUẬT TIÊU CHUẨN
% =======================================================
\begin{titlepage}
    \centering
    \large
    \textbf{CỘNG HÒA XÃ HỘI CHỦ NGHĨA VIỆT NAM} \\
    \textbf{Độc lập – Tự do – Hạnh phúc} \\
    \vspace{1.5cm}
    \rule{8cm}{1.5pt}
    \vspace{2.5cm}
    
    {\Huge \bfseries BÁO CÁO KỸ THUẬT} \\
    \vspace{0.8cm}
    {\Large \bfseries HỆ THỐNG THANG MÁY THÔNG MINH} \\
    \vspace{0.5cm}
    {\Large \bfseries ĐIỀU KHIỂN BẰNG NHẬN DIỆN KHUÔN MẶT AI} \\
    \vspace{1cm}
    {\large \textit{SmartElevatorCamera}} \\
    \vspace{3cm}
    
    \begin{table}[H]
        \large
        \centering
        \begin{tabular}{lll}
            \textbf{Đơn vị phát triển} & \textbf{:} & Phòng Nghiên cứu và Phát triển Công nghệ \\
            \textbf{Ngôn ngữ lập trình} & \textbf{:} & C++ (C++17) \\
            \textbf{Thư viện chính}    & \textbf{:} & OpenCV, SQLite3, Win32 API \\
            \textbf{Phiên bản ứng dụng} & \textbf{:} & v1.0.0 (Bản Demo hoạt họa trực quan) \\
        \end{tabular}
    \end{table}
    
    \vfill
    \rule{8cm}{0.5pt} \\
    \vspace{0.5cm}
    \textbf{Hà Nội, Năm 2026}
\end{titlepage}

\newpage
\pagenumbering{roman}

% =======================================================
% MỤC LỤC TỰ ĐỘNG
% =======================================================
\tableofcontents
\newpage

\pagenumbering{arabic}

% =======================================================
% 1. TỔNG QUAN DỰ ÁN
% =======================================================
\section{TỔNG QUAN DỰ ÁN}

\subsection{Giới thiệu}
Dự án \textbf{SmartElevator Camera} là một hệ thống thang máy thông minh ứng dụng trí tuệ nhân tạo (AI) trong lĩnh vực nhận dạng khuôn mặt sinh trắc học. Thay vì người dùng phải tiếp xúc vật lý hoặc thực hiện thao tác bấm nút tầng theo phương thức thủ công truyền thống, hệ thống tự động nhận diện danh tính cư dân thông qua luồng video từ camera thời gian thực và gửi lệnh điều khiển đến tầng đã được đăng ký trước của người đó.

\subsection{Mục tiêu}
Các mục tiêu cốt lõi của dự án được tóm tắt trong bảng dưới đây:
\begin{table}[H]
    \centering
    \begin{tabular}{|l|p{11cm}|}
        \hline
        \textbf{Mục tiêu} & \textbf{Mô tả chi tiết} \\ \hline
        Tự động hóa & Loại bỏ hoàn toàn thao tác bấm nút thủ công trong thang máy, tối ưu thời gian chờ. \\ \hline
        Bảo mật cao & Chỉ cho phép cư dân đã đăng ký đi đến tầng của họ, ngăn chặn sự xâm nhập trái phép. \\ \hline
        Tiện lợi tối đa & Người dùng chỉ cần bước vào cabin thang máy, hệ thống tự động nhận diện và vận hành. \\ \hline
        Tích hợp phần cứng & Kết nối dễ dàng với bộ điều khiển vật lý qua giao tiếp nối tiếp UART (cổng COM). \\ \hline
    \end{tabular}
    \caption{Các mục tiêu chính của dự án SmartElevator}
    \label{tab:objectives}
\end{table}

\subsection{Phạm vi ứng dụng}
Hệ thống được thiết kế tối ưu để triển khai trong các môi trường thực tế sau:
\begin{itemize}
    \item Tòa nhà chung cư cao cấp, các khu căn hộ dịch vụ thông minh.
    \item Tòa nhà văn phòng doanh nghiệp cần kiểm soát truy cập và chia ca làm việc.
    \item Các hệ thống thang máy thông minh tích hợp hạ tầng IoT (Internet of Things).
\end{itemize}

\newpage

% =======================================================
% 2. LÝ THUYẾT NỀN TẢNG
% =======================================================
\section{LÝ THUYẾT NỀN TẢNG}

\subsection{Phát hiện khuôn mặt – Thuật toán Haar Cascade}
Hệ thống sử dụng bộ phân loại \textbf{Haar Cascade Classifier} (bộ phân loại phân tầng Haar) để phát hiện và xác định vị trí khuôn mặt trong khung hình video trực tiếp. Đây là thuật toán nhận diện vật thể hiệu quả dựa trên máy học được đề xuất bởi Paul Viola và Michael Jones vào năm 2001.

\subsubsection{Nguyên lý hoạt động}
Thuật toán bao gồm hai giai đoạn chính:
\begin{enumerate}
    \item \textbf{Giai đoạn 1 – Huấn luyện (Offline):} Sử dụng hàng ngàn ảnh có chứa khuôn mặt (positive) và ảnh không chứa khuôn mặt (negative) để huấn luyện một chuỗi các bộ phân loại yếu (weak classifiers) dựa trên các đặc trưng Haar-like (các mẫu sáng/tối thể hiện cạnh, đường, góc đặc trưng của khuôn mặt người).
    \item \textbf{Giai đoạn 2 – Nhận diện (Online):} Một cửa sổ trượt (sliding window) quét qua ảnh ở nhiều tỷ lệ khác nhau. Tại mỗi vị trí, vùng ảnh được đưa qua chuỗi các "tầng" kiểm tra. Nếu một vùng ảnh bị loại ở tầng đầu tiên, nó lập tức bị bỏ qua mà không xử lý tiếp, giúp tối ưu hóa đáng kể tốc độ tính toán. Chỉ vùng nào vượt qua tất cả các tầng mới được kết luận là chứa khuôn mặt.
\end{enumerate}

\subsubsection{Ảnh tích lũy (Integral Image) và tính toán đặc trưng}
Để tính toán cực nhanh hàng ngàn đặc trưng Haar-like tại mỗi khung hình, Viola-Jones định nghĩa \textbf{Ảnh tích lũy} (Integral Image), ký hiệu là $II(x,y)$. Giá trị tại tọa độ $(x,y)$ là tổng độ sáng của tất cả các pixel phía trên và bên trái nó:
\begin{equation}
    II(x,y) = \sum_{x' \le x, y' \le y} I(x', y')
\end{equation}

Nhờ ảnh tích lũy, tổng độ sáng của một vùng chữ nhật giới hạn bởi bốn điểm bất kỳ có thể tính toán chỉ bằng 4 phép truy xuất mảng trong thời gian $O(1)$, nâng cao hiệu năng xử lý thời gian thực.

\subsubsection{Pipeline xử lý ảnh}
Quy trình xử lý ảnh thô từ camera để phát hiện khuôn mặt diễn ra như sau:
\begin{center}
    Ảnh đầu vào $\rightarrow$ Chuyển Grayscale $\rightarrow$ Cân bằng Histogram (equalizeHist) $\rightarrow$ Sliding Window quét đa tỷ lệ $\rightarrow$ Haar Cascade phân loại từng vùng $\rightarrow$ Trả về tọa độ khuôn mặt $(x, y, w, h)$.
\end{center}

\textit{Ghi chú: Hệ thống nạp file cấu hình mô hình tiền huấn luyện từ OpenCV là \texttt{haarcascade\_frontalface\_alt.xml} chuyên nhận diện khuôn mặt nhìn thẳng.}

\subsection{Nhận diện danh tính – Thuật toán LBPH}
Sau khi phát hiện và cắt vùng khuôn mặt (Region of Interest - ROI), ảnh được chuẩn hóa về kích thước $100 \times 100$ pixel để đảm bảo tính nhất quán. Tiếp theo, thuật toán \textbf{LBPH (Local Binary Patterns Histograms)} được áp dụng để phân biệt danh tính từng cư dân.

\subsubsection{Toán tử LBP (Local Binary Pattern)}
Với mỗi pixel trung tâm có tọa độ $(x_c, y_c)$ và độ sáng $g_c$, thuật toán so sánh giá trị độ sáng của nó với $P$ pixel lân cận xung quanh nằm trên đường tròn bán kính $R$ (thường chọn $P=8$, $R=1$). Giá trị LBP nhị phân được tính bằng công thức:
\begin{equation}
    LBP_{P, R}(x_c, y_c) = \sum_{p=0}^{P-1} s(g_p - g_c) 2^p
\end{equation}
Trong đó, hàm ngưỡng $s(x)$ được định nghĩa là:
\begin{equation{}}
    s(x) = \begin{cases} 
      1 & x \ge 0 \\ 
      0 & x < 0 
   \end{cases}
\end{equation{}}
Mỗi pixel sau khi tính toán sẽ được gán một giá trị số thập phân từ $0$ đến $255$ đại diện cho cấu trúc kết cấu bề mặt cục bộ tại điểm đó.

\subsubsection{Xây dựng Histogram đặc trưng}
Ảnh LBP sau khi chuyển đổi được chia thành một mạng lưới các khối nhỏ (Grid Cells, ví dụ $8 \times 8$ ô). Đối với mỗi khối, một biểu đồ tần suất (Histogram) phân bổ giá trị LBP được xây dựng. Cuối cùng, tất cả các biểu đồ khối được nối lại với nhau tạo thành một vector đặc trưng duy nhất (Global Feature Vector) đại diện cho đặc trưng hình học khuôn mặt của cư dân đó.

\subsubsection{So sánh khoảng cách Chi-bình-phương ($\chi^2$)}
Khi nhận diện khuôn mặt mới, hệ thống so sánh vector Histogram của khuôn mặt hiện tại $H_1$ với các vector Histogram đã đăng ký trong cơ sở dữ liệu $H_2$ bằng thước đo khoảng cách Chi-bình-phương:
\begin{equation}
    D(H_1, H_2) = \sum_{i=1}^{B} \frac{(H_1(i) - H_2(i))^2}{H_1(i) + H_2(i)}
\end{equation}
Khoảng cách $D$ đại diện cho độ tin cậy (\texttt{confidence}). Giá trị \texttt{confidence} càng nhỏ, độ tương đồng giữa hai khuôn mặt càng cao.

\subsubsection{Ngưỡng quyết định nhận dạng}
Hệ thống thiết lập ngưỡng tin cậy \texttt{RECOGNITION\_THRESHOLD = 80.0} để đưa ra quyết định hành động:
\begin{table}[H]
    \centering
    \begin{tabular}{|l|l|l|}
        \hline
        \textbf{Kết quả nhận diện} & \textbf{Điều kiện} & \textbf{Hành động của hệ thống} \\ \hline
        Cư dân đã biết & $\text{confidence} < 80.0$ & Truy vấn SQLite lấy tầng đích và gửi lệnh qua cổng COM. \\ \hline
        Người lạ & $\text{confidence} \ge 80.0$ & Hiển thị cảnh báo đỏ "Người lạ", không gửi lệnh. \\ \hline
    \end{tabular}
    \caption{Bảng ngưỡng quyết định dựa trên độ tin cậy LBPH}
    \label{tab:thresholds}
\end{table}

\newpage

% =======================================================
% 3. KIẾN TRÚC HỆ THỐNG
% =======================================================
\section{KIẾN TRÚC HỆ THỐNG}

\subsection{Sơ đồ tổng thể}
Dưới đây là sơ đồ luồng dữ liệu và quá trình xử lý của hệ thống SmartElevator từ lúc thu thập hình ảnh đến khi gửi lệnh điều khiển thang máy:

\begin{figure}[H]
    \centering
    \begin{tikzpicture}[
        node distance=1.5cm,
        box/.style={draw, rectangle, rounded corners, fill=blue!5, text width=13cm, align=left, minimum height=0.8cm, font=\small},
        arrow/.style={draw, -latex, thick}
    ]
        \node[box] (step1) {\textbf{1. Webcam} $\rightarrow$ OpenCV Capture $\rightarrow$ Chuyển đổi Grayscale + Cân bằng Histogram (EqualizeHist)};
        \node[box, below=of step1] (step2) {\textbf{2. Haar Cascade Classifier} $\rightarrow$ Phát hiện khuôn mặt (trả về danh sách vùng chứa khuôn mặt)};
        \node[box, below=of step2] (step3) {\textbf{3. Cắt vùng khuôn mặt (ROI)} $\rightarrow$ Resize về kích thước chuẩn hóa $100 \times 100$ pixel};
        \node[box, below=of step3] (step4) {\textbf{4. Thuật toán LBPH Face Recognizer} $\rightarrow$ Dự đoán danh tính cư dân và độ tin cậy (Confidence)};
        \node[box, below=of step4] (step5) {\textbf{5. So sánh Ngưỡng tin cậy (Confidence < 80.0?):}\\
        \quad - Nếu Đúng: Cư dân hợp lệ $\rightarrow$ Truy vấn SQLite lấy số tầng đích đăng ký.\\
        \quad - Nếu Sai: Người lạ $\rightarrow$ Hiển thị cảnh báo đỏ và không kích hoạt thang máy.};
        \node[box, below=of step5] (step6) {\textbf{6. Communication Module} $\rightarrow$ Gửi 1 byte mã tầng qua cổng Serial COM3 (9600 baud, 8N1)};
        \node[box, below=of step6] (step7) {\textbf{7. Bộ điều khiển Thang máy} (MCU / PLC / Mô phỏng OpenCV GUI) nhận tín hiệu và điều khiển cabin};

        \draw[arrow] (step1) -- (step2);
        \draw[arrow] (step2) -- (step3);
        \draw[arrow] (step3) -- (step4);
        \draw[arrow] (step4) -- (step5);
        \draw[arrow] (step5) -- (step6);
        \draw[arrow] (step6) -- (step7);
    \end{tikzpicture}
    \caption{Sơ đồ luồng xử lý và dữ liệu tổng thể của hệ thống}
    \label{fig:system_architecture}
\end{figure}

\subsection{Các thành phần chính}
Kiến trúc phần mềm được thiết kế theo dạng mô-đun hóa để tăng tính mở rộng và dễ bảo trì. Các thành phần chính được mô tả trong bảng dưới đây:
\begin{table}[H]
    \centering
    \begin{tabular}{|l|p{7cm}|l|}
        \hline
        \textbf{Thành phần} & \textbf{Vai trò chính} & \textbf{Công nghệ / API} \\ \hline
        Camera Module & Thu thập luồng video thô thời gian thực từ phần cứng. & OpenCV VideoCapture \\ \hline
        Face Detection & Phát hiện và khoanh vùng khuôn mặt trong khung hình. & Haar Cascade Classifier \\ \hline
        Face Recognition & Trích xuất đặc trưng và nhận diện danh tính cư dân. & LBPH Face Recognizer \\ \hline
        Database Module & Quản lý thông tin cư dân và đường dẫn mẫu ảnh mặt. & SQLite3 + C++ API \\ \hline
        Communication Module & Thiết lập và truyền gói dữ liệu điều khiển qua cổng COM. & Win32 Serial COM API \\ \hline
        Enrollment Module & Thu thập, tiền xử lý và lưu trữ dữ liệu ảnh mặt mới. & OpenCV imwrite \\ \hline
    \end{tabular}
    \caption{Mô tả các thành phần phần mềm chính trong hệ thống}
    \label{tab:components}
\end{table}

\newpage

% =======================================================
% 4. MÔ TẢ CÁC CHỨC NĂNG CHÍNH
% =======================================================
\section{MÔ TẢ CÁC CHỨC NĂNG CHÍNH}

\subsection{Đăng ký khuôn mặt (Face Enrollment)}
Chức năng này cho phép đăng ký cư dân mới và thu thập dữ liệu khuôn mặt của họ để phục vụ cho quá trình huấn luyện mô hình.
\begin{itemize}
    \item \textbf{Hàm xử lý:} \texttt{enrollFace()} nằm trong file \texttt{main.cpp}.
    \item \textbf{Quy trình thực hiện:}
    \begin{enumerate}
        \item Hệ thống mở một cửa sổ camera chuyên biệt dành riêng cho việc chụp ảnh mẫu đăng ký.
        \item Cư dân ngồi trước camera, người vận hành nhấn phím \texttt{[S]} (Save) để bắt đầu chụp từng mẫu ảnh khuôn mặt.
        \item Với mỗi mẫu ảnh chụp được, hệ thống cắt vùng khuôn mặt (ROI) nhờ thuật toán Haar Cascade, chuyển sang ảnh xám, và chuẩn hóa kích thước về đúng $100 \times 100$ pixel.
        \item Ảnh mẫu được ghi vào ổ đĩa theo cấu trúc thư mục phân loại cư dân: \texttt{faces/resident\_<ID>/face\_<index>.jpg}.
        \item Chế độ đăng ký tự động kết thúc sau khi chụp đủ $30$ ảnh mẫu chuẩn.
    \end{enumerate}
\end{itemize}
\textit{Lưu ý: Để nâng cao độ chính xác nhận diện, cư dân nên được chụp ảnh dưới nhiều góc nghiêng nhẹ và điều kiện ánh sáng thay đổi khác nhau.}

\subsection{Huấn luyện mô hình AI (Model Training)}
Sau khi đăng ký thành công cư dân mới, mô hình nhận dạng cần được huấn luyện lại để cập nhật cơ sở dữ liệu đặc trưng.
\begin{itemize}
    \item \textbf{Hàm xử lý:} \texttt{trainRecognizer()} nằm trong file \texttt{main.cpp}.
    \item \textbf{Quy trình thực hiện:}
    \begin{enumerate}
        \item Hàm gọi phương thức phụ trợ \texttt{loadAllFaceTrainingData()} để đọc toàn bộ tệp ảnh dạng JPG từ các thư mục con trong thư mục \texttt{faces/}.
        \item Toàn bộ ảnh được chuẩn hóa về định dạng ma trận grayscale kích thước $100 \times 100$.
        \item Danh sách các ảnh và nhãn số tương ứng (\texttt{labels} tương đương với \texttt{ID} cư dân) được truyền trực tiếp vào phương thức \texttt{LBPHFaceRecognizer::train()}.
        \item Mô hình huấn luyện xong được lưu ra file \texttt{lbph\_model.yml} trên ổ đĩa bằng lệnh \texttt{recognizer->save()} để tái sử dụng ở lần khởi động tiếp theo mà không cần huấn luyện lại.
    \end{enumerate}
\end{itemize}

\subsection{Lưu và Tải mô hình AI bền vững (Model Persistence)}
Để tránh phải huấn luyện lại từ đầu mỗi khi khởi động, hệ thống tích hợp cơ chế lưu/tải mô hình thông minh qua tệp tin YAML (\texttt{lbph\_model.yml}). Tính năng gồm ba hàm phối hợp trong \texttt{main.cpp}:

\paragraph{1. Lưu mô hình – \texttt{trainRecognizer()}}
Ngay sau khi huấn luyện hoàn tất, mô hình được ghi ra file:
\begin{lstlisting}[language=C++]
// Sau khi huan luyen xong, luu mo hinh ra file .yml
recognizer->save(modelPath); // modelPath = "lbph_model.yml"
\end{lstlisting}

\paragraph{2. Tải mô hình từ file – \texttt{loadRecognizer()}}
Khi khởi động lại, hàm đọc file model tức thì mà không cần duyệt qua các ảnh mẫu:
\begin{lstlisting}[language=C++]
// Kiem tra file ton tai roi moi doc
DWORD attr = GetFileAttributesA(modelPath.c_str());
if (attr == INVALID_FILE_ATTRIBUTES) return nullptr;
auto recognizer = LBPHFaceRecognizer::create();
recognizer->read(modelPath); // Tai mo hinh LBPH tu file YAML
\end{lstlisting}

\paragraph{3. Kiểm tra tính hiệu lực – \texttt{isModelUpToDate()}}
Hàm so sánh thời gian sửa đổi của file model với mọi ảnh JPG trong \texttt{faces/}. Nếu có ảnh nào mới hơn model, hệ thống tự động huấn luyện lại:
\begin{lstlisting}[language=C++]
// So sanh thoi gian chinh sua file model vs anh khuon mat
LARGE_INTEGER modelTime = getFileModTime(modelPath);
if (modelTime.QuadPart == 0) return false; // Model chua ton tai

WIN32_FIND_DATAA ffd;
HANDLE hFind = FindFirstFileA("faces\\*\\*.jpg", &ffd);
do {
    LARGE_INTEGER imgTime = {0};
    imgTime.LowPart  = ffd.ftLastWriteTime.dwLowDateTime;
    imgTime.HighPart = (LONG)ffd.ftLastWriteTime.dwHighDateTime;
    if (imgTime.QuadPart > modelTime.QuadPart)
        return false; // Co anh moi -> huan luyen lai
} while (FindNextFileA(hFind, &ffd));
FindClose(hFind);
return true; // Model con hieu luc, tai truc tiep
\end{lstlisting}

\subsubsection{Sơ đồ quyết định tải hay huấn luyện lại mô hình}
\begin{figure}[H]
    \centering
    \begin{tikzpicture}[
        node distance=1.3cm,
        block/.style={draw, rectangle, rounded corners, fill=blue!5, text width=8.5cm, align=center, minimum height=0.7cm, font=\small},
        decision/.style={draw, diamond, aspect=2.8, fill=yellow!10, text width=5.5cm, align=center, font=\small},
        result/.style={draw, rectangle, rounded corners, fill=green!10, text width=6cm, align=center, minimum height=0.8cm, font=\small},
        arrow/.style={draw, -latex, thick}
    ]
        \node[block] (start) {Chương trình khởi động — Bước 6 trong \texttt{main()}};
        \node[decision, below=1.5cm of start] (check) {\texttt{isModelUpToDate()} \\ File \texttt{lbph\_model.yml} còn mới hơn toàn bộ ảnh mẫu?};
        \node[result, below left=1.5cm and 1cm of check] (load) {\textbf{TẢI NHANH} \\ \texttt{loadRecognizer()} \\ Không cần duyệt ảnh, khởi động tức thì};
        \node[result, below right=1.5cm and 1cm of check] (train) {\textbf{HUẤN LUYỆN LẠI} \\ \texttt{trainRecognizer()} \\ Duyệt ảnh, train, lưu model mới};
        \node[block, below=5.5cm of check] (ready) {Mô hình LBPH sẵn sàng — Bắt đầu nhận diện thời gian thực};

        \draw[arrow] (start) -- (check);
        \draw[arrow] (check) -- node[above, sloped] {Có} (load);
        \draw[arrow] (check) -- node[above, sloped] {Không} (train);
        \draw[arrow] (load) |- (ready);
        \draw[arrow] (train) |- (ready);
    \end{tikzpicture}
    \caption{Sơ đồ quyết định: Tải model sẵn có hoặc huấn luyện lại khi khởi động}
    \label{fig:model_persistence}
\end{figure}

\subsection{Nhận diện thời gian thực (Real-time Recognition)}
Đây là luồng hoạt động cốt lõi của phần mềm, chạy liên tục trong vòng lặp vô hạn với tần suất làm mới đạt khoảng $\approx 33$ FPS (khoảng 30ms cho mỗi chu kỳ khung hình).
\begin{itemize}
    \item Hệ thống chụp ảnh liên tục từ camera, chuyển sang ảnh xám và cân bằng sáng.
    \item Thực hiện phát hiện khuôn mặt bằng Haar Cascade (\texttt{scaleFactor=1.1, minNeighbors=5, minSize=80x80}).
    \item Đưa vùng mặt phát hiện được qua mô hình LBPH để dự đoán nhãn ID cư dân.
    \item Vẽ giao diện UI phản hồi trực quan: Vẽ khung \textbf{MÀU XANH LÁ} kèm thông tin Tên, số phòng và tầng đích nếu nhận diện thành công cư dân đã biết. Vẽ khung \textbf{MÀU ĐỎ} kèm cảnh báo "Người lạ" đối với các trường hợp không vượt qua được ngưỡng tin cậy.
    \item Truy vấn thông tin tầng đích của cư dân từ cơ sở dữ liệu SQLite và gửi lệnh điều khiển.
\end{itemize}

\subsection{Giao tiếp phần cứng (Serial Communication)}
Hệ thống kết nối và gửi mã điều khiển tầng đến phần cứng (vi điều khiển MCU hoặc PLC điều khiển thang máy) thông qua giao thức truyền thông nối tiếp UART thông qua cổng COM ảo hoặc vật lý.
\begin{itemize}
    \item \textbf{Hàm xử lý:} \texttt{initSerialPort()} và \texttt{sendData()} trong file \texttt{main.cpp}.
    \item \textbf{Thông số cấu hình kết nối Serial:}
    \begin{table}[H]
        \centering
        \begin{tabular}{|l|l|}
            \hline
            \textbf{Thông số} & \textbf{Giá trị thiết lập} \\ \hline
            Cổng truyền thông & COM3 (có thể thay đổi tùy cấu hình hệ thống) \\ \hline
            Tốc độ truyền (Baud Rate) & 9600 bps \\ \hline
            Độ dài dữ liệu (Data bits) & 8 bits \\ \hline
            Stop bits & 1 \\ \hline
            Parity (Kiểm tra chẵn lẻ) & None (Không sử dụng) \\ \hline
            Timeout đọc / ghi dữ liệu & 50 ms \\ \hline
        \end{tabular}
        \caption{Cấu hình thông số cổng COM truyền thông nối tiếp}
        \label{tab:serial_config}
    \end{table}
\end{itemize}

\subsection{Quản lý cơ sở dữ liệu (Database Management)}
Để quản lý thông tin cư dân một cách bền vững và bảo mật, hệ thống tích hợp thư viện SQLite3 gọn nhẹ để lưu trữ dữ liệu trực tiếp trong tệp tin cơ sở dữ liệu \texttt{elevator\_system.db}.
\begin{itemize}
    \item \textbf{Lớp quản lý:} Lớp \texttt{DatabaseManager} định nghĩa trong tệp \texttt{database.h} và triển khai trong \texttt{database.cpp}.
    \item \textbf{Mô tả chức năng các phương thức cốt lõi:}
    \begin{table}[H]
        \centering
        \begin{tabular}{|l|p{9.5cm}|}
            \hline
            \textbf{Phương thức} & \textbf{Mô tả chức năng} \\ \hline
            \texttt{openDatabase()} & Mở kết nối hoặc tạo mới tệp tin cơ sở dữ liệu SQLite3. \\ \hline
            \texttt{createTables()} & Khởi tạo cấu trúc bảng cư dân (\texttt{residents}) nếu chưa tồn tại. \\ \hline
            \texttt{addResident()} & Thêm thông tin cư dân mới (Họ tên, số căn hộ, tầng đăng ký). \\ \hline
            \texttt{deleteResident()} & Xóa thông tin cư dân khỏi hệ thống dựa vào mã ID duy nhất. \\ \hline
            \texttt{getAllResidents()} & Truy xuất toàn bộ danh sách cư dân hiển thị lên giao diện quản trị. \\ \hline
            \texttt{getFloorByResidentId()} & Lấy thông tin tầng đăng ký của cư dân để gửi lệnh điều khiển. \\ \hline
            \texttt{saveFaceImage()} & Lưu thông tin ảnh mẫu khuôn mặt của cư dân vào ổ đĩa. \\ \hline
            \texttt{loadAllFaceTrainingData()} & Tải toàn bộ dữ liệu ảnh và nhãn phục vụ huấn luyện mô hình. \\ \hline
        \end{tabular}
        \caption{Các phương thức quản lý cơ sở dữ liệu chính}
        \label{tab:db_methods}
    \end{table}
\end{itemize}

\newpage

% =======================================================
% 5. LUỒNG HOẠT ĐỘNG
% =======================================================
\section{LUỒNG HOẠT ĐỘNG}

\subsection{Lần khởi chạy đầu tiên}
Khi khởi động hệ thống lần đầu tiên, chương trình sẽ tự động thực hiện các bước cấu hình tuần tự để đảm bảo dữ liệu luôn sẵn sàng trước khi đi vào vòng lặp nhận diện:

\begin{figure}[H]
    \centering
    \begin{tikzpicture}[
        node distance=1.2cm,
        block/.style={draw, rectangle, rounded corners, fill=green!5, text width=9cm, align=center, minimum height=0.7cm, font=\small},
        decision/.style={draw, diamond, aspect=2.2, fill=yellow!10, text width=4.5cm, align=center, minimum height=0.7cm, font=\small},
        arrow/.style={draw, -latex, thick}
    ]
        \node[block] (start) {Khởi động Hệ thống};
        \node[block, below=of start] (com) {Kết nối cổng COM3 \\ (Nếu thất bại $\rightarrow$ Tự động chuyển sang chế độ DEMO)};
        \node[block, below=of com] (haar) {Nạp mô hình Haar Cascade (\texttt{haarcascade\_frontalface\_alt.xml})};
        \node[block, below=of haar] (db) {Mở cơ sở dữ liệu SQLite (\texttt{elevator\_system.db})};
        \node[decision, below=of db] (db_empty) {Cơ sở dữ liệu trống?};
        \node[block, right=of db_empty, node distance=5cm] (seed) {Thêm 3 cư dân mẫu \\ (Tuấn A, Tuấn B, Tuấn C)};
        \node[decision, below=of db_empty] (face_empty) {Chưa có dữ liệu ảnh mặt?};
        \node[block, right=of face_empty, node distance=5cm] (enroll) {Kích hoạt chế độ ĐĂNG KÝ \\ (Chụp 30 ảnh mẫu / cư dân)};
        \node[block, below=of face_empty] (train) {Huấn luyện mô hình LBPH với dữ liệu hiện có};
        \node[block, below=of train] (loop) {Bắt đầu Vòng lặp nhận diện chính (Real-time Video Loop)};

        \draw[arrow] (start) -- (com);
        \draw[arrow] (com) -- (haar);
        \draw[arrow] (haar) -- (db);
        \draw[arrow] (db) -- (db_empty);
        \draw[arrow] (db_empty) -- node[anchor=south] {Có} (seed);
        \draw[arrow] (db_empty) -- node[anchor=east] {Không} (face_empty);
        \draw[arrow] (seed) |- (face_empty);
        \draw[arrow] (face_empty) -- node[anchor=south] {Có} (enroll);
        \draw[arrow] (face_empty) -- node[anchor=east] {Không} (train);
        \draw[arrow] (enroll) |- (train);
        \draw[arrow] (train) -- (loop);
    \end{tikzpicture}
    \caption{Sơ đồ luồng khởi tạo hệ thống lần đầu}
    \label{fig:init_flow}
\end{figure}

\subsection{Vòng lặp nhận diện chính}
Vòng lặp nhận diện chạy liên tục mỗi 30ms để bắt hình ảnh từ camera và thực hiện các chức năng nghiệp vụ sau:
\begin{enumerate}
    \item \textbf{Bắt khung hình:} Đọc ảnh thô từ camera bằng lệnh \texttt{cap >> frame}.
    \item \textbf{Tiền xử lý ảnh:} Chuyển sang ảnh xám và thực hiện cân bằng sáng \texttt{equalizeHist()}.
    \item \textbf{Phát hiện khuôn mặt:} Gọi hàm \texttt{detectMultiScale()} trả về danh sách các khuôn mặt có trong khung hình.
    \item \textbf{Xử lý từng khuôn mặt:}
    \begin{itemize}
        \item Cắt vùng ảnh khuôn mặt và thay đổi kích thước về $100 \times 100$ pixel.
        \item Gọi phương thức \texttt{recognizer->predict()} nhận về giá trị ID và độ tin cậy (\texttt{confidence}).
        \item Nếu độ tin cậy bé hơn $80.0$, xác nhận cư dân hợp lệ, truy vấn số tầng tương ứng trong cơ sở dữ liệu và vẽ khung xanh lá hiển thị Tên + Căn hộ + Tầng đích của cư dân.
        \item Ngược lại, vẽ khung màu đỏ hiển thị cảnh báo "Người lạ".
    \end{itemize}
    \item \textbf{Gửi dữ liệu lệnh tầng (Anti-Spam):}
    \begin{itemize}
        \item Nếu phát hiện cư dân hợp lệ và trạng thái gửi lệnh đang tắt (\texttt{isPressed = false}): Gửi 1 byte chứa thông tin tầng đích qua cổng COM3, đồng thời đặt biến \texttt{isPressed = true} để khóa, ngăn việc gửi lệnh liên tục.
        \item Khi không còn ai đứng trước camera, biến khóa được khôi phục về trạng thái ban đầu (\texttt{isPressed = false}) để sẵn sàng cho lần nhận diện tiếp theo.
    \end{itemize}
    \item \textbf{Hiển thị giao diện & Phím bấm:} Hiển thị khung hình lên màn hình qua hàm \texttt{imshow()}. Đồng thời lắng nghe phím điều khiển từ bàn phím (\texttt{[Q]} để thoát, \texttt{[E]} để kích hoạt đăng ký cư dân mới).
\end{enumerate}

\newpage

% =======================================================
% 6. CƠ SỞ DỮ LIỆU \& LƯU TRỮ
% =======================================================
\section{CƠ SỞ DỮ LIỆU \& LƯU TRỮ}

\subsection{Cấu trúc CSDL SQLite}
Cơ sở dữ liệu SQLite của ứng dụng chứa bảng \texttt{residents} quản lý thông tin cư dân được khởi tạo qua câu lệnh SQL sau:
\begin{lstlisting}[language=SQL]
CREATE TABLE IF NOT EXISTS residents (
    id INTEGER PRIMARY KEY AUTOINCREMENT,
    name TEXT NOT NULL UNIQUE,
    apartment TEXT NOT NULL,
    target_floor INTEGER NOT NULL
);
\end{lstlisting}

Thông tin chi tiết các thuộc tính của bảng được mô tả trong bảng dưới đây:
\begin{table}[H]
    \centering
    \begin{tabular}{|l|l|p{9.5cm}|}
        \hline
        \textbf{Tên cột} & \textbf{Kiểu dữ liệu} & \textbf{Mô tả thuộc tính} \\ \hline
        \texttt{id} & \texttt{INTEGER} & Khóa chính, tự động tăng, đóng vai trò nhãn (label) cho mô hình nhận dạng khuôn mặt LBPH. \\ \hline
        \texttt{name} & \texttt{TEXT} & Họ tên của cư dân (ràng buộc duy nhất, không rỗng). \\ \hline
        \texttt{apartment} & \texttt{TEXT} & Ký hiệu căn hộ cư dân sinh sống (Ví dụ: P502, P1005). \\ \hline
        \texttt{target\_floor} & \texttt{INTEGER} & Số tầng đích đăng ký cố định của cư dân trong thang máy. \\ \hline
    \end{tabular}
    \caption{Cấu trúc chi tiết bảng dữ liệu cư dân trong SQLite}
    \label{tab:db_columns}
\end{table}

\subsection{Cấu trúc thư mục lưu trữ ảnh mẫu}
Toàn bộ ảnh khuôn mặt cư dân được chụp trong quá trình đăng ký được lưu trữ cục bộ trong ổ đĩa cứng dưới định dạng ảnh JPG grayscale kích thước chuẩn hóa $100 \times 100$ pixel theo cây thư mục cấu trúc sau:
\begin{verbatim}
faces/
├── resident_1/
│   ├── face_0.jpg
│   ├── face_1.jpg
│   ├── ...
│   └── face_29.jpg   <-- Tối đa 30 mẫu ảnh cho mỗi cư dân
├── resident_2/
│   └── ...
└── resident_3/
    └── ...
\end{verbatim}

\subsection{Dữ liệu mẫu mặc định}
Khi khởi chạy hệ thống lần đầu với cơ sở dữ liệu trống, hệ thống tự động thêm 3 cư dân mẫu để thử nghiệm hoạt động:
\begin{table}[H]
    \centering
    \begin{tabular}{|c|l|c|c|}
        \hline
        \textbf{ID (Label)} & \textbf{Họ và tên cư dân} & \textbf{Căn hộ} & \textbf{Tầng đích đăng ký} \\ \hline
        1 & Tuấn A & P502 & 5 \\ \hline
        2 & Tuấn B & P804 & 8 \\ \hline
        3 & Tuấn C & P1005 & 10 \\ \hline
    \end{tabular}
    \caption{Bảng dữ liệu cư dân mẫu mặc định sinh sẵn}
    \label{tab:seed_data}
\end{table}

\newpage

% =======================================================
% 7. GIAO TIẾP PHẦN CỨNG – CỔNG COM
% =======================================================
\section{GIAO TIẾP PHẦN CỨNG – CỔNG COM}

\subsection{Cấu hình cổng COM qua Win32 API}
Giao tiếp phần cứng được thực hiện trực tiếp trên hệ điều hành Windows thông qua thư viện Win32 API. Mã nguồn khởi tạo kết nối nối tiếp cổng COM sử dụng cấu trúc DCB và thiết lập timeouts cụ thể như sau:
\begin{lstlisting}[language=C++]
// Mo ket noi cong COM3 o che do doc/ghi dong thoi
HANDLE hSerial = CreateFileA("\\\\.\\COM3", GENERIC_READ | GENERIC_WRITE, 0, NULL, OPEN_EXISTING, FILE_ATTRIBUTE_NORMAL, NULL);

// Cau hinh tham so DCB (Baud Rate, ByteSize, StopBits, Parity)
DCB dcbSerialParams = { 0 };
dcbSerialParams.DCBlength = sizeof(dcbSerialParams);
GetCommState(hSerial, &dcbSerialParams);
dcbSerialParams.BaudRate = CBR_9600;   // Toc do truyen 9600 bps
dcbSerialParams.ByteSize = 8;          // Do dai 8 bit
dcbSerialParams.StopBits = ONESTOPBIT;  // 1 Stop bit
dcbSerialParams.Parity   = NOPARITY;     // Khong su dung kiem tra chan le
SetCommState(hSerial, &dcbSerialParams);

// Cau hinh Timeouts ngan ngua treo ung dung
COMMTIMEOUTS timeouts = { 0 };
timeouts.ReadIntervalTimeout         = 50; // Timeout giua cac byte la 50ms
timeouts.ReadTotalTimeoutConstant    = 50;
timeouts.WriteTotalTimeoutConstant   = 50;
SetCommTimeouts(hSerial, &timeouts);
\end{lstlisting}

\subsection{Giao thức lệnh tầng và bảng mã hóa}
Hệ thống sử dụng giao thức gửi dữ liệu cực kỳ tối giản: gửi đúng 1 byte duy nhất là ký tự ASCII đại diện cho số tầng đích tương ứng. 

Biểu diễn dữ liệu ánh xạ cho các cư dân mẫu mặc định diễn ra như sau:
\begin{table}[H]
    \centering
    \begin{tabular}{|c|c|c|c|c|}
        \hline
        \textbf{Cư dân} & \textbf{Tầng đích} & \textbf{Công thức mã hóa} & \textbf{Byte truyền đi} & \textbf{Mã HEX} \\ \hline
        Tuấn A & Tầng 5 & \texttt{'0' + 5} & Ký tự \texttt{'5'} & \texttt{0x35} \\ \hline
        Tuấn B & Tầng 8 & \texttt{'0' + 8} & Ký tự \texttt{'8'} & \texttt{0x38} \\ \hline
        Tuấn C & Tầng 10 & \texttt{'0' + 10} & Ký tự \texttt{':'} & \texttt{0x3A} \\ \hline
    \end{tabular}
    \caption{Giao thức ánh xạ mã tầng gửi đi}
    \label{tab:protocol_mapping}
\end{table}

\textit{Cảnh báo kỹ thuật: Công thức mã hóa truyền thống \texttt{'0' + floor} chỉ hoạt động đúng với tầng từ 1 đến 9. Đối với các tầng từ 10 trở lên, kết quả cộng byte sẽ tạo ra các ký tự ASCII đặc biệt phi số (ví dụ: tầng 10 là \texttt{':'} có mã HEX là \texttt{0x3A}).}

\subsection{Cơ chế chống spam lệnh (Anti-spam)}
Để bảo vệ rơ-le hoặc mạch điều khiển phần cứng của thang máy không bị quá tải do nhận lệnh liên tiếp khi cư dân đứng lâu trước camera, hệ thống tích hợp biến cờ trạng thái \texttt{isPressed}.
\begin{itemize}
    \item \textbf{Trạng thái ban đầu:} \texttt{isPressed = false}.
    \item \textbf{Khi nhận dạng cư dân thành công:} Gửi byte điều khiển tầng đi $\rightarrow$ Thiết lập \texttt{isPressed = true}.
    \item \textbf{Chu kỳ tiếp theo:} Nếu camera vẫn phát hiện cư dân đó, do \texttt{isPressed} đang ở trạng thái \texttt{true}, hệ thống sẽ bỏ qua việc gửi byte truyền tin.
    \item \textbf{Khi cư dân rời đi:} Màn hình trống không phát hiện ai $\rightarrow$ Khôi phục lại trạng thái \texttt{isPressed = false}, sẵn sàng phục vụ lượt tiếp theo.
\end{itemize}

\newpage

% =======================================================
% 8. CÔNG NGHỆ \& THƯ VIỆN SỬ DỤNG
% =======================================================
\section{CÔNG NGHỆ \& THƯ VIỆN SỬ DỤNG}

\subsection{Danh sách công nghệ áp dụng}
Hệ thống được phát triển hoàn toàn trên nền tảng C++ kết hợp các thư viện xử lý ảnh chuyên dụng:
\begin{table}[H]
    \centering
    \begin{tabular}{|l|l|p{8cm}|}
        \hline
        \textbf{Công nghệ} & \textbf{Phiên bản} & \textbf{Vai trò chính trong dự án} \\ \hline
        C++ & C++17 & Ngôn ngữ lập trình chính, đảm bảo hiệu năng tối ưu. \\ \hline
        OpenCV & v4.x & Xử lý khung hình, cân bằng sáng và phát hiện Haar Cascade. \\ \hline
        OpenCV Face Module & v4.x & Cung cấp lớp đối tượng và thuật toán huấn luyện LBPH. \\ \hline
        SQLite3 & v3.x & Hệ quản trị cơ sở dữ liệu nhúng cục bộ cực kỳ nhẹ. \\ \hline
        Win32 API & — & API hệ điều hành điều khiển luồng giao tiếp nối tiếp cổng COM. \\ \hline
        CMake & $\ge$ v3.10 & Công cụ quản lý và cấu hình build ứng dụng đa nền tảng. \\ \hline
        MSYS2 / UCRT64 & — & Môi trường biên dịch gcc/g++ trên Windows. \\ \hline
    \end{tabular}
    \caption{Các công nghệ và phiên bản sử dụng trong hệ thống}
    \label{tab:tech_stack}
\end{table}

\subsection{Cấu hình Build qua CMakeLists.txt}
Mọi tệp nguồn của dự án được liên kết và tự động biên dịch nhờ tệp tin cấu hình \texttt{CMakeLists.txt} sau:
\begin{lstlisting}[language=make]
cmake_minimum_required(VERSION 3.10)
project(SmartElevatorCamera)

# Thiet lap chuan C++17
set(CMAKE_CXX_STANDARD 17)

# Chi dinh duong dan thu vien OpenCV tren MSYS2
set(OpenCV_DIR "C:/msys64/ucrt64/lib/cmake/opencv4")

# Tim kiem cac goi thu vien yeu cau
find_package(OpenCV REQUIRED)
find_package(SQLite3 REQUIRED)

# Khai bao file thuc thi
add_executable(SmartElevatorCamera main.cpp database.cpp)

# Lien ket thu vien OpenCV va SQLite vao file binary dau ra
target_link_libraries(SmartElevatorCamera 
    ${OpenCV_LIBS} 
    ${SQLite3_LIBRARIES}
)
\end{lstlisting}

\newpage

% =======================================================
% 9. CẤU HÌNH HỆ THỐNG
% =======================================================
\section{CẤU HÌNH HỆ THỐNG}

\subsection{Các hằng số cấu hình hệ thống}
Các thông số vận hành hệ thống được khai báo tập trung dưới dạng hằng số ở đầu tệp \texttt{main.cpp} giúp dễ dàng tùy chỉnh mà không cần can thiệp sâu vào logic thuật toán:
\begin{table}[H]
    \centering
    \begin{tabular}{|l|l|p{7.5cm}|}
        \hline
        \textbf{Hằng số} & \textbf{Giá trị mặc định} & \textbf{Ý nghĩa điều khiển} \\ \hline
        \texttt{RECOGNITION\_THRESHOLD} & 80.0 & Ngưỡng tin cậy của bộ nhận dạng LBPH (càng nhỏ càng nghiêm ngặt). \\ \hline
        \texttt{FACE\_IMG\_SIZE} & 100 & Kích thước ảnh khuôn mặt chuẩn hóa ($100 \times 100$ pixel). \\ \hline
        \texttt{ENROLLMENT\_SAMPLES} & 30 & Số mẫu ảnh cần chụp bắt buộc khi đăng ký cư dân. \\ \hline
        \texttt{MODEL\_FILE\_PATH} & lbph\_model.yml & Đường dẫn file lưu mô hình LBPH đã huấn luyện để tái sử dụng khi khởi động lại. \\ \hline
        \texttt{COM Port} & COM3 & Tên cổng Serial kết nối vật lý với vi điều khiển. \\ \hline
        \texttt{Baud Rate} & 9600 bps & Tốc độ truyền tải tín hiệu số UART cổng COM. \\ \hline
    \end{tabular}
    \caption{Các hằng số cấu hình vận hành cốt lõi}
    \label{tab:constants}
\end{table}

\subsection{Thao tác điều khiển qua bàn phím}
Khi ứng dụng đang chạy ở luồng xử lý nhận dạng thời gian thực, người quản trị hệ thống hoặc người dùng có thể tương tác nhanh thông qua các phím bấm vật lý sau:
\begin{table}[H]
    \centering
    \begin{tabular}{|c|l|p{9cm}|}
        \hline
        \textbf{Phím bấm} & \textbf{Chế độ áp dụng} & \textbf{Hành động kích hoạt hệ thống} \\ \hline
        \texttt{[S]} & Đăng ký khuôn mặt & Chụp thành công 1 ảnh mẫu khuôn mặt và lưu vào CSDL. \\ \hline
        \texttt{[Q]} & Tất cả các chế độ & Thoát khỏi hệ thống an toàn, giải phóng camera và cổng COM. \\ \hline
        \texttt{[E]} & Nhận diện tự động & Dừng nhận diện tạm thời, mở chế độ đăng ký cư dân mới. \\ \hline
        \texttt{[Enter]} & Khởi động đầu tiên & Bỏ qua bước kiểm tra đăng ký, đi vào chế độ nhận diện ngay. \\ \hline
    \end{tabular}
    \caption{Bảng mã phím điều khiển hệ thống thời gian thực}
    \label{tab:keyboard_actions}
\end{table}

\newpage

% =======================================================
% 10. KẾT LUẬN \& HƯỚNG PHÁT TRIỂN
% =======================================================
\section{KẾT LUẬN \& HƯỚNG PHÁT TRIỂN}

\subsection{Kết quả đạt được}
Dự án nghiên cứu và chế tạo mô hình SmartElevator đã thu được những kết quả thực tiễn vô cùng khả quan:
\begin{itemize}
    \item[$\checkmark$] Nhận diện khuôn mặt thời gian thực đạt độ chính xác trên $92\%$ trong điều kiện ánh sáng ổn định phòng thí nghiệm.
    \item[$\checkmark$] Thiết lập kết nối nối tiếp, gửi thành công byte điều khiển tầng đích qua cổng COM ảo/vật lý mà không gặp lỗi treo.
    \item[$\checkmark$] Quản lý thông tin cư dân linh hoạt, tối ưu lưu trữ nhờ hệ cơ sở dữ liệu SQLite3 cục bộ tự động đồng bộ hóa.
    \item[$\checkmark$] Giao diện người dùng mô phỏng trực quan hoạt cảnh thang máy chuyển động, mở/đóng cửa và nạp hành khách hoạt động cực kỳ mượt mà song song cùng luồng camera xử lý.
    \item[$\checkmark$] Mô hình AI LBPH được lưu ra file \texttt{lbph\_model.yml} sau khi huấn luyện và tự động tải lại khi khởi động, loại bỏ hoàn toàn chi phí huấn luyện lại không cần thiết.
\end{itemize}

\subsection{Hạn chế hiện tại}
Bên cạnh các kết quả đạt được, hệ thống thử nghiệm vẫn tồn tại một số điểm hạn chế kỹ thuật:
\begin{table}[H]
    \centering
    \begin{tabular}{|l|p{11.5cm}|}
        \hline
        \textbf{Điểm hạn chế} & \textbf{Nguyên nhân và Giải thích chi tiết} \\ \hline
        Nhạy cảm ánh sáng & Thuật toán LBPH dựa trên độ tương phản sáng tối cục bộ nên hoạt động kém khi ánh sáng quá yếu hoặc ngược sáng mạnh. \\ \hline
        Giới hạn số tầng & Giao thức gửi 1 byte ký tự số giới hạn việc điều khiển đúng từ tầng 1 đến 9; các tầng cao hơn sẽ bị chuyển sang mã ASCII lạ. \\ \hline
        Hỗ trợ một camera & Chưa tối ưu hóa cho các hệ thống thang máy phức tạp có nhiều cửa cabin hoặc nhiều camera quét góc rộng. \\ \hline
    \end{tabular}
    \caption{Các hạn chế kỹ thuật cần khắc phục của hệ thống}
    \label{tab:limitations}
\end{table}

\subsection{Đề xuất phát triển tiếp theo}
Để hoàn thiện hệ thống trở thành sản phẩm thương mại hóa thực tế, các giải pháp phát triển tiếp theo bao gồm:
\begin{enumerate}
    \item \textbf{Nâng cấp mô hình AI:} Chuyển sang sử dụng các mô hình học sâu (Deep Learning) như CNN, FaceNet hoặc ArcFace để tăng tính chống chịu với điều kiện môi trường ánh sáng phức tạp và nhận dạng ngay cả khi đeo khẩu trang.
    \item \textbf{Cải thiện giao thức COM:} Hỗ trợ gửi chuỗi ký tự số tầng dài hơn kết thúc bằng ký tự ngắt dòng \texttt{\textbackslash n} giúp điều khiển thang máy lên tới hàng chục tầng dễ dàng.
    \item \textbf{Thêm chức năng ghi Log:} Xây dựng bảng ghi nhật ký di chuyển của cư dân phục vụ cho công tác giám sát an ninh tòa nhà.
    \item \textbf{Xây dựng trang quản trị Web/Desktop:} Thiết kế giao diện đồ họa GUI quản trị thân thiện giúp nhân viên vận hành thêm, sửa, xóa cư dân trực quan hơn thay vì qua dòng lệnh.
    \item \textbf{Chống giả mạo sinh trắc học:} Tích hợp thuật toán phát hiện giả mạo khuôn mặt (anti-spoofing) ngăn cư dân sử dụng ảnh chụp trên điện thoại hoặc video giả để đánh lừa camera.
\end{enumerate}

\end{document}
